\label{chap.7}
\section{Why MLOps}
A forecasting model trained once and left alone slowly goes stale: new weather arrives every day, the climate shifts (rainfall in all five districts has risen over the 26 years, Chapter~\ref{chap.3}), and nobody notices when accuracy drops. Sculley et al.\ \cite{sculley2015hidden} point out that in real machine-learning systems the model is a small part of the whole, and most of the effort goes into data handling, configuration, serving and monitoring. MLOps is the practice of treating those parts as carefully as the model. For this project it means four things: the data and models are versioned, there is one clearly marked production model per district, forecasts are served through a stable interface, and performance is checked automatically.

\section{Architecture}
Figure~\ref{fig:mlops} shows the flow. The pipeline is built around the project's main forecasting model, the per-district hurdle rainfall model; the comparative models are kept as research notebooks and are not served.

\begin{figure}[H]
\centering
\begin{tikzpicture}[
  font=\sffamily\small,
  box/.style={draw, rounded corners, fill=blue!6, minimum width=3.6cm, minimum height=0.9cm, align=center},
  arr/.style={-{Stealth[length=2.2mm]}, thick}
]
\node[box] (n) {NASA POWER API};
\node[box, below=0.55cm of n] (d) {Weather data\\(versioned with DVC)};
\node[box, below=0.55cm of d] (m) {Hurdle model per district};
\node[box, below=0.55cm of m] (r) {MLflow registry\\(\emph{champion} version)};
\node[box, below=0.55cm of r] (a) {FastAPI forecast service};
\node[box, below=0.55cm of a] (l) {Forecast log};
\node[box, right=1.8cm of l] (mo) {Weekly monitor};
\node[box, right=1.8cm of m] (c) {Challenger model};
\draw[arr] (n) -- node[right, font=\sffamily\footnotesize]{weekly refresh} (d);
\draw[arr] (d) -- (m); \draw[arr] (m) -- (r); \draw[arr] (r) -- (a); \draw[arr] (a) -- (l);
\draw[arr] (l) -- (mo);
\draw[arr] (mo) -- node[right, font=\sffamily\footnotesize, align=left]{accuracy\\dropped?} (c);
\draw[arr] (c) |- node[pos=0.75, above, font=\sffamily\footnotesize]{promote if better} (r);
\end{tikzpicture}
\caption{MLOps flow for the forecasting stage}
\label{fig:mlops}
\end{figure}

\section{Components}
\begin{table}[H]
	\centering
	\caption{MLOps components}
	\label{tab:mlops}
	\small
	\begin{tabularx}{\textwidth}{p{3.4cm}p{4.2cm}Y}
		\toprule
		\textbf{Component} & \textbf{File / tool} & \textbf{What it does} \\
		\midrule
		Configuration & \texttt{params.yaml} & One place for districts, data paths, model settings (horizon 3, thresholds 0.25 mm and 0.35), split dates and monitoring thresholds \\
		Data refresh & \texttt{fetch\_nasa.py}, \texttt{clean.py} & Download the latest 60 days from NASA POWER and clean them into the weather folder \\
		Data versioning & DVC with a DagsHub remote & Keeps every version of the data and model outputs outside Git, with a small pointer file in Git \\
		Model pipeline & \texttt{hurdle\_pipeline.py} & Feature engineering, leakage-safe per-horizon training and evaluation, forecasting \\
		Registry & \texttt{registry.py}, MLflow & Logs each model with its parameters and metrics, versions it, and marks one version per district as the \emph{champion} \\
		Serving & \texttt{serve/app.py}, FastAPI, Docker & \texttt{/health} and \texttt{/forecast/\{district\}} endpoints; every forecast is appended to a log \\
		Monitoring & \texttt{monitor.py} & Scores logged forecasts against observed rain and, if accuracy has dropped, trains and tests a challenger \\
		Automation & GitHub Actions & CI tests on every push; weekly data refresh; weekly monitoring \\
		Tests & \texttt{pytest} & Check feature engineering and the hurdle model contract \\
		\bottomrule
	\end{tabularx}
\end{table}

\section{How it works}

\subsection{Data versioning}
The weather files are too large and change too often to keep in Git. DVC (Data Version Control) stores them in a remote store on DagsHub and keeps only a small pointer file in Git. Every model can therefore be traced to the exact data it was trained on, and any earlier state can be restored with \texttt{dvc pull}.

\subsection{Model registry and the champion}
MLflow \cite{zaharia2018accelerating} records every trained model together with its parameters (district, horizon, test start date, probability threshold) and test metrics. Each district has one registered model, \texttt{<district>\_hurdle}, with numbered versions. One version carries the alias \emph{champion}; the API always serves the champion, so switching to a better model means moving the alias rather than editing code.

\subsection{Serving}
The FastAPI service exposes \texttt{GET /health} and \texttt{GET /forecast/\{district\}}. On a request it fetches the district's champion from MLflow (and caches it), loads the latest weather, returns the forecast as JSON, and appends a line to \texttt{forecasts\_log.csv} with the model version, the date the forecast was made, the target date and the predicted rainfall.

\subsection{Automation}
Three GitHub Actions workflows run without manual effort:
\begin{itemize}
	\item \textbf{ci}: on every push or pull request, install the dependencies and run the tests;
	\item \textbf{weekly-data}: every Monday at 03:00 UTC (08:30 IST), pull the current data, fetch new NASA POWER data, clean it, push the new version to DVC and commit the updated pointer;
	\item \textbf{weekly-monitor}: every Monday at 06:00 UTC (11:30 IST), score the logged forecasts.
\end{itemize}

\subsection{Monitoring and champion--challenger promotion}
Once at least 30 days of logged forecasts can be compared with observed rainfall, the monitor computes RMSE and $R^2$ over the last 60 days. If $R^2$ has fallen more than 0.10 below the champion's recorded test $R^2$, it trains a \emph{challenger} on the latest data, evaluates both models on the same most recent year, and promotes the challenger only if it is better. This keeps a worse model from replacing a good one by accident.

\section{Current status and honest limitations}
\begin{itemize}
	\item The pipeline, registry, API, workflows and tests are in place and the five existing hurdle checkpoints have been registered.
	\item The registered checkpoints are the \emph{original} hurdle models, from the first, same-day version of the model (Section~\ref{sec:leak}). The API therefore labels its output \texttt{original\_checkpoint\_same\_day}. Registering the corrected three-horizon models described in Chapter~\ref{chap.5}, already implemented in \texttt{hurdle\_pipeline.py}, is the first task of the next phase.
	\item The monitor needs the forecast log to be versioned with DVC before the scheduled workflow can read it; this is a manual step today.
	\item The comparative models (LSTM, SARIMA, XGBoost, TFT) are not registered or served; if a later stage needs forecasts of other variables such as temperature, the LSTM can be added to the same pipeline.
\end{itemize}
